\documentclass[10pt,letterpaper]{article}

\usepackage[margin=1in]{geometry}
\usepackage{amsmath,amssymb,amsfonts}
\usepackage{booktabs}
\usepackage{multirow}
\usepackage{graphicx}
\usepackage{xcolor}
\usepackage{hyperref}
\usepackage{enumitem}
\usepackage{microtype}
\usepackage[numbers]{natbib}
\usepackage{caption}

\definecolor{civo}{RGB}{180,40,40}
\definecolor{healthy}{RGB}{40,120,40}
\definecolor{partial}{RGB}{180,140,30}

\title{Diagnosing and Mitigating Value Overestimation\\from Temporal Abstraction in DreamerV3}

\author{Anonymous Authors}

\date{}

\begin{document}

\maketitle

%% ============================================================
%%  ABSTRACT
%% ============================================================
\begin{abstract}
Temporal abstraction---executing actions at variable rates---is a long-sought capability in reinforcement learning, yet integrating it into world-model agents remains poorly understood.
We present a systematic diagnostic study of variable-rate imagination in DreamerV3, identifying a failure mode we term \emph{conditioning-induced value overestimation} (CIVO): when a trained temporal-step embedding (\texttt{dt\_embedding}) conditions the world model, the critic systematically overestimates values for imagined states at $dt>1$, triggering a positive feedback loop that collapses the policy.
Through factorial experiments spanning 4 environments (Walker Walk, Cheetah Run, Cartpole Swingup, Hopper Hop), 30+ configurations, and 60+ seeds, we establish that
(i)~a normalized value-to-evaluation ratio (NVE) $\geq 1.5\times$ at 35K+ steps reliably diagnoses CIVO (AUROC\,=\,0.969, 26/26 classification across 4 environments),
(ii)~CIVO arises from forward-pass representational damage, not gradient contamination---all backward-pass (selective gradient stopping) interventions fail (1/17 seeds, 3 variants, 2 environments),
and (iii)~a forward-pass intervention, \emph{Temporal Conditioning Noise} (TCN), which injects Gaussian noise into the \texttt{dt\_embedding} during imagination, successfully prevents CIVO while retaining 90\% baseline performance (Walker, $\sigma=0.2$) and reduces NVE from 8--108$\times$ to 1.5--4$\times$ (Cheetah, $\sigma=0.5$).
A dose-response relationship between untreated CIVO severity and required noise magnitude provides additional mechanistic insight.
All claims are scoped to DreamerV3 and its RSSM world model.
\end{abstract}

%% ============================================================
%%  1. INTRODUCTION
%% ============================================================
\section{Introduction}

Temporal abstraction---the ability to reason and act at multiple time scales---is a foundational desideratum in reinforcement learning \cite{sutton1999between,precup2000temporal}.
In model-based RL, world models offer a natural substrate: by conditioning imagination on a temporal step size $dt$, an agent can simulate future trajectories at variable rates without modifying environment interaction.
Recent successes with DreamerV3 \cite{hafner2023mastering} have made latent-dynamics world models the de facto platform for such experiments.

Yet a persistent puzzle emerges: every attempt to equip DreamerV3 with variable-rate imagination---whether via learned gating, fixed $dt>1$, or temporal subsampling---degrades asymptotic performance, often catastrophically.
This stands in contrast to the success of temporal conditioning in model-predictive control (MPC) settings, where TAWM \cite{nhu2025tawm} demonstrates effective variable-rate planning with TD-MPC2.
Why does temporal abstraction succeed in MPC but fail in imagination-based actor-critic?

This paper provides a systematic, causally controlled answer.
Through factorial experiments spanning 4 DeepMind Control Suite environments, 30+ configurations, and 60+ seeds, we identify the root cause: a failure mode we term \textbf{conditioning-induced value overestimation (CIVO)}.
When a trained \texttt{dt\_embedding} conditions the RSSM world model, the critic exploits precise but inaccurate temporal conditioning to systematically overestimate values at $dt>1$, triggering a positive feedback loop (overestimation $\to$ actor exploitation $\to$ distribution shift $\to$ worse predictions $\to$ greater overestimation) that collapses the policy.

Our investigation proceeds in three stages:
\begin{enumerate}[leftmargin=*,nosep]
    \item \textbf{Diagnosis} (C1--C5): We isolate trained \texttt{dt\_embedding} as the necessary and sufficient condition for CIVO through controlled ablations, establishing a diagnostic threshold (NVE\,$\geq$\,1.5$\times$) with near-perfect classification across environments.
    \item \textbf{Mechanism} (C6--C7): We show CIVO operates through forward-pass representational damage, not gradient contamination. All three backward-pass interventions (selective gradient stopping) fail across 17 seeds and 2 environments, while the trained embedding creates vulnerability through at least two distinct pathways.
    \item \textbf{Mitigation} (C8): Guided by the forward-pass diagnosis, we propose \emph{Temporal Conditioning Noise} (TCN), which injects calibrated Gaussian noise into the \texttt{dt\_embedding} during imagination. TCN prevents CIVO in Walker ($\sigma=0.2$, NVE\,$<$\,1.0$\times$) and substantially reduces it in the more severe Cheetah environment ($\sigma=0.5$, NVE from 8--108$\times$ to 1.5--4$\times$), with a dose-response relationship linking noise magnitude to untreated severity.
\end{enumerate}

All claims are scoped to DreamerV3 and its RSSM world model.
Generalizability to other latent-dynamics architectures is discussed but not claimed.

%% ============================================================
%%  2. BACKGROUND & RELATED WORK
%% ============================================================
\section{Background and Related Work}

\paragraph{Temporal abstraction in RL.}
The options framework \cite{sutton1999between,precup2000temporal} formalizes temporal abstraction as policies over variable-length sub-routines.
FiGAR \cite{sharma2017learning} learns action repetition counts.
In model-free settings, simple action repeat \cite{lakshminarayanan2017dynamic} is the most common form.

\paragraph{World-model RL.}
The Dreamer family \cite{hafner2020dream,hafner2021mastering,hafner2023mastering} learns latent-dynamics models (RSSM) for imagination-based policy optimization.
IRIS \cite{micheli2023transformers} and TWM \cite{robine2023transformer} use transformer world models.
These methods train actor and critic entirely within imagined trajectories.

\paragraph{Variable-rate planning in world models.}
TAWM \cite{nhu2025tawm} successfully conditions TD-MPC2's world model on temporal step size for variable-rate MPC planning.
The key architectural difference is that MPC re-plans from real environment states at each step, while imagination-based actor-critic bootstraps from imagined states---creating the positive feedback loop that enables CIVO.

\paragraph{Value overestimation in MBRL.}
Model exploitation---where the policy exploits inaccuracies in the learned model---is a known challenge \cite{janner2019trust}.
Lu et al.~\cite{lu2018non} identify ``delusional bias'' from policy-class constraints.
Our CIVO mechanism is distinct: it arises from \emph{conditioning-induced} exploitation of temporal embedding inaccuracies, not from policy constraints or general model error.

%% ============================================================
%%  3. METHODOLOGY
%% ============================================================
\section{Methodology}

\subsection{DreamerV3 with Variable-Rate Imagination}

We extend DreamerV3 \cite{hafner2023mastering} with temporal abstraction components:
\begin{itemize}[leftmargin=*,nosep]
    \item \textbf{Temporal subsampling}: Training data includes transitions at $dt=1$ and $dt>1$ (obtained by skipping intermediate frames).
    \item \textbf{dt\_embedding}: A learned embedding $e_{dt} \in \mathbb{R}^d$ conditions the RSSM, reward head, and observation decoder on the temporal step size.
    \item \textbf{Variable-rate imagination}: During actor-critic training, imagined rollouts can use $dt>1$ steps.
\end{itemize}

\subsection{Factorial Experiment Design}

We use a controlled factorial design varying three factors:
\begin{enumerate}[leftmargin=*,nosep]
    \item \textbf{dt\_embedding state}: trained, untrained/unexposed (present but never sees $dt>1$ data), frozen (random fixed), or absent ($\text{dim}=0$).
    \item \textbf{Temporal subsampling}: ON (training data contains $dt>1$ transitions) or OFF.
    \item \textbf{Imagination horizon}: $H=8$ or $H=15$; imagination dt: $dt=1$ or $dt=2$.
\end{enumerate}
Experiments span 4 DMC environments: Walker Walk, Cheetah Run, Cartpole Swingup, and Hopper Hop.

\subsection{Diagnostic Metric: Normalized Value-to-Evaluation Ratio (NVE)}

We define the \emph{normalized value-to-evaluation ratio}:
\begin{equation}
    \text{NVE}(t) = \frac{V(t) / E(t)}{V_{\text{base}}(t) / E_{\text{base}}(t)},
\end{equation}
where $V(t)$ is the mean predicted value, $E(t)$ is the mean evaluation return, and subscript ``base'' denotes the $dt=1$ baseline at the same training step.
NVE\,$\geq$\,1.5$\times$ at $t \geq 35$K indicates pathological overestimation; NVE\,$<$\,1.0$\times$ indicates value suppression.
The 35K lower bound ensures the baseline value-to-evaluation ratio has stabilized (at 25K, baseline $V/E$ is only 0.087, causing ratio instability).

%% ============================================================
%%  4. RESULTS
%% ============================================================
\section{Experiments and Results}

All experiments use DeepMind Control Suite \cite{tassa2018deepmind} environments.
The DreamerV3 $dt=1$ baseline achieves evaluation returns of $\sim$916 (Walker Walk), $\sim$594 (Cheetah Run), $\sim$785 (Cartpole Swingup), and $\sim$213 (Hopper Hop), averaged across seeds.

%% ---- 4.1 C1 ----
\subsection{Variable-Rate Imagination Consistently Degrades Performance (C1)}

Every $dt>1$ imagination configuration degrades asymptotic performance relative to the $dt=1$ baseline.
Fixed $dt=2$ with trained embedding achieves peak 485 before collapsing to $\sim$135.
Fixed $dt=4$ reaches only 177 (33\% of baseline at the same step).
A learned gating mechanism peaks at 541 before collapsing to 99.
The degradation is monotonic in $dt$: $dt=1$ ($\sim$950) $>$ $dt=2$ (peak 485) $>$ $dt=4$ (peak 177).
This holds across all tested environments.

%% ---- 4.2 C2 ----
\subsection{Trained dt\_embedding Reliably Triggers CIVO (C2)}
\label{sec:c2}

The central diagnostic finding: trained \texttt{dt\_embedding} is the necessary and sufficient condition for CIVO.

\paragraph{Walker Walk cross-seed validation (008d, trained embedding).}
Both seeds exceed NVE\,$\geq$\,1.5$\times$ at 35K+ steps.
Seed~1 sustains elevated NVE across three consecutive checkpoints (1.71$\times$\,$\to$\,2.09$\times$\,$\to$\,4.12$\times$) and collapses (eval 131 at 105K).
Seed~2 transiently exceeds 1.5$\times$ at 45K (NVE\,=\,1.62$\times$) but self-corrects (NVE\,=\,0.38$\times$ at 105K, eval 268).

\begin{table}[h]
\centering
\caption{Walker 008d (trained dt\_embedding, subsample ON): cross-seed NVE trajectory.}
\label{tab:008d_cross_seed}
\small
\begin{tabular}{@{}rcccc@{}}
\toprule
Step & Seed 1 NVE & Seed 1 Eval & Seed 2 NVE & Seed 2 Eval \\
\midrule
25K  & 1.78 & 167.2 & 2.05 & 141.1 \\
35K  & 1.71 & 146.3 & 1.24 & 164.1 \\
45K  & 2.09 & 140.8 & 1.62 & 141.0 \\
55K  & \textbf{4.12} & 82.7  & 1.41 & 172.3 \\
65K  & 1.23 & 250.0 & 0.94 & 209.2 \\
105K & 1.57 & 131.1 & 0.38 & 267.9 \\
\bottomrule
\end{tabular}
\end{table}

\paragraph{Causal control: no dt\_embedding (008f).}
Removing the embedding entirely ($\text{dt\_embed\_dim}=0$) eliminates CIVO across 4 seeds: all NVE\,$<$\,1.5$\times$ (range 0.85--1.41).
Performance degrades (148--503, vs.\ baseline $\sim$916) due to reward prediction collapse (reward MAE\,=\,0.747, 75$\times$ baseline), but no seed collapses via CIVO.

\begin{table}[h]
\centering
\caption{Walker 008f (dt\_embed\_dim=0): all seeds show NVE\,$<$\,1.5$\times$, no CIVO.}
\label{tab:008f}
\small
\begin{tabular}{@{}lcccc@{}}
\toprule
Seed & Eval@105K & NVE max (35K+) & Collapse & CIVO \\
\midrule
s0 & 503 & 0.97 & No & No \\
s1 & 380 & 0.85 & No & No \\
s2 & 148 & 1.14 & No & No \\
s3 & 364 & $<$1.5 & No & No \\
s4 & 307 & 1.41 & No & No \\
\bottomrule
\end{tabular}
\end{table}

\paragraph{Frozen embedding (015b).}
Freezing the \texttt{dt\_embedding} at random initialization ($\text{requires\_grad}=\text{False}$) produces a third failure mode---gradient starvation (NVE\,$=$\,0.42, actor loss$\to$0)---distinct from CIVO.
This confirms that \emph{training} of the embedding, not its mere presence, is the necessary condition.

\begin{table}[h]
\centering
\caption{dt\_embedding state $\times$ dt$>$1 data presence: failure mode matrix.}
\label{tab:2x2_matrix}
\small
\begin{tabular}{@{}lll@{}}
\toprule
dt\_emb state & $dt>1$ data present & $dt>1$ data absent \\
\midrule
Trained  & \textcolor{civo}{CIVO collapse (NVE 2--4$\times$)} & Value suppression (NVE 0.5--0.7$\times$) \\
Frozen   & Gradient starvation (NVE 0.4$\times$) & n/a \\
Absent   & No CIVO, reward degraded (NVE 0.3--1.1$\times$) & $=$ baseline \\
\bottomrule
\end{tabular}
\end{table}

\paragraph{RSSM fidelity gap.}
The mechanistic basis of CIVO is a fidelity gap in the RSSM's predictions at $dt>1$.
Table~\ref{tab:rssm_fidelity} shows that the trained embedding (008d) amplifies the reward prediction gap to 2.78$\times$ ($dt=2$ vs.\ $dt=1$ MAE), creating the exploitation channel.
Removing the embedding (008f) collapses this gap to 1.01$\times$ (at the cost of uniformly poor reward prediction, MAE\,=\,0.75).

\begin{table}[h]
\centering
\caption{RSSM prediction fidelity: absolute errors and $dt=2$/$dt=1$ ratios. Both absolute values and ratios are reported; 008d's low obs gap (1.30$\times$) is misleading because its $dt=1$ error (41.3) is already 3$\times$ baseline---the model has collapsed.}
\label{tab:rssm_fidelity}
\small
\begin{tabular}{@{}lcccccc@{}}
\toprule
Config & \multicolumn{3}{c}{Obs MSE} & \multicolumn{3}{c}{Reward MAE} \\
\cmidrule(lr){2-4}\cmidrule(lr){5-7}
 & $dt$=1 & $dt$=2 & Gap & $dt$=1 & $dt$=2 & Gap \\
\midrule
Baseline          & 13.49 & 27.50 & 2.04$\times$ & 0.010 & 0.016 & 1.56$\times$ \\
008d (trained)    & 41.31 & 53.67 & 1.30$\times$ & 0.057 & 0.158 & \textbf{2.78$\times$} \\
008f (no emb)     & 8.01  & 22.04 & 2.75$\times$ & 0.747 & 0.753 & 1.01$\times$ \\
\bottomrule
\end{tabular}
\end{table}

%% ---- 4.3 C3 ----
\subsection{Untrained dt\_embedding Prevents Optimistic Drift (C3)}

When the \texttt{dt\_embedding} is present but never exposed to $dt>1$ data (subsample OFF), it cannot learn meaningful temporal discrimination.
All 3 Walker seeds show NVE\,$<$\,1.0$\times$ (range 0.51--0.66) with performance at 85\% of baseline (eval range 631--899).
This pattern replicates in Cheetah (3 seeds, NVE 0.48--0.50, 87\% baseline).

\begin{table}[h]
\centering
\caption{Unexposed dt\_embedding: cross-environment value suppression.}
\label{tab:unexposed_cross_env}
\small
\begin{tabular}{@{}lcccc@{}}
\toprule
Environment & Eval vs.\ Baseline & NVE Range & Seeds $<$1.0$\times$ \\
\midrule
Walker Walk  & 85\% (778/916) & 0.51--0.66 & 3/3 \\
Cheetah Run  & 87\% (519/594) & 0.48--0.50 & 3/3 \\
\bottomrule
\end{tabular}
\end{table}

%% ---- 4.4 C4+C5 ----
\subsection{Temporal Step Size and Horizon Effects (C4, C5)}

CIVO severity is monotonic in $dt$: $dt=2$ (peak eval 485, collapse to 135) is strictly worse than $dt=1$ ($\sim$950) but better than $dt=4$ (peak 177, 33\% of baseline).
This is consistent with RSSM prediction error compounding over longer temporal gaps.

Subsampling (shorter effective training horizons) contributes to degradation but is neither necessary nor sufficient for collapse.
With subsample OFF and untrained embedding, $H=15$ achieves 85\% baseline (778) while $H=8$ achieves only 26\% (242), indicating horizon length matters independently.
Critically, with subsample ON and trained embedding, restricting imagination to $dt=1$ does not prevent collapse (Section~\ref{sec:c6}), confirming that the root cause is embedding-mediated overestimation, not data reduction.

%% ---- 4.5 C6 ----
\subsection{Dual Collapse Pathways from Trained dt\_embedding (C6)}
\label{sec:c6}

A key finding: restricting imagination to $dt=1$ while retaining the trained \texttt{dt\_embedding} does \emph{not} prevent collapse.
Across 4 independent seeds, 2/4 (50\%) collapse---the same rate as $dt=2$ imagination.

\begin{table}[h]
\centering
\caption{008e\_v3 (trained dt\_emb, $dt=1$ imagination): 4-seed results. Collapse rate 50\%.}
\label{tab:008e_v3_seeds}
\small
\begin{tabular}{@{}lcccl@{}}
\toprule
Seed & Eval@105K & NVE Peak & Collapse & Pathway \\
\midrule
s1 & 869.9 & 0.80 & No  & --- \\
s2 & 137.0 & 3.78 & Yes & B (eval-driven) \\
s3 & 92.7  & 4.02 & Yes & B (eval-driven) \\
s4 & 899.3 & 0.99 & No  & --- \\
\bottomrule
\end{tabular}
\end{table}

The collapsed seeds exhibit dynamics \emph{qualitatively different} from classic CIVO (Pathway~A).
We identify two distinct pathways:

\begin{table}[h]
\centering
\caption{Two collapse pathways: classic CIVO (Pathway A) vs.\ policy instability (Pathway B).}
\label{tab:dual_pathway}
\small
\begin{tabular}{@{}lccc@{}}
\toprule
Feature & Seed 1 (healthy) & 008e\_v3 s2 (Pathway B) & 008d s0 (Pathway A) \\
\midrule
Eval@105K            & 869.9              & 137.0                & 82.7 \\
NVE@105K             & 0.8                & 3.78                 & 4.12 \\
Value $\Delta$ 85K$\to$105K & +30\%       & +0.1\%               & +130\% \\
Eval $\Delta$ 85K$\to$105K  & +32\%       & $-$81\%              & $-$44\% \\
NVE driver           & n/a                & Eval denominator collapse & Value numerator explosion \\
Entropy trend        & Stable ($-$5.2)    & $\uparrow$ ($-$5.2$\to$$-$3.8)  & $\downarrow$ ($-$3.5$\to$$-$4.9) \\
\bottomrule
\end{tabular}
\end{table}

\textbf{Pathway~A} (classic CIVO): The critic overestimates $\to$ the actor over-exploits high-valued actions $\to$ entropy decreases $\to$ policy becomes deterministic and wrong.
\textbf{Pathway~B} (policy instability): The RSSM's distorted latent representations (from $dt>1$ training data, even without $dt>1$ imagination) cause inconsistent actor gradients $\to$ entropy increases $\to$ policy degrades toward random.

A baseline run with the same seed (009a, seed~2) is completely stable (eval 618$\to$596$\to$837), ruling out seed-specific fragility.
The trained \texttt{dt\_embedding} creates vulnerability through multiple pathways; $dt=1$ imagination eliminates Pathway~A but not Pathway~B.

Actor entropy serves as a discriminator between the three regimes:

\begin{table}[h]
\centering
\caption{Actor entropy trajectories distinguish three regimes.}
\label{tab:entropy}
\small
\begin{tabular}{@{}llccc@{}}
\toprule
Trajectory Type & Representative & Early & Late & Direction \\
\midrule
Pathway A (CIVO)       & 008d s0           & $-$3.47 & $-$4.88 & $\downarrow$ exploitation \\
Pathway B (instability) & 008e\_v3 s2      & $-$5.21 & $-$3.83 & $\uparrow$ randomization \\
Healthy                 & baseline / 008a   & $-$4.50 & $-$5.10 & $\to$ stable \\
\bottomrule
\end{tabular}
\end{table}

%% ---- 4.6 C7 ----
\subsection{Backward-Pass Interventions Fail (C7)}
\label{sec:c7}

We hypothesized that selectively stopping gradients from specific loss components through the \texttt{dt\_embedding} (selective gradient stopping, SGS) could prevent CIVO by blocking the co-adaptation channel.
Three SGS variants were tested across 17 seeds and 2 environments:

\begin{enumerate}[leftmargin=*,nosep]
    \item \textbf{Reward-only SGS}: Detach \texttt{dt\_embedding} from reward loss gradients. Despite the reward head contributing the largest gradient fraction (38--49\%), this intervention \emph{worsens} CIVO (NVE\,=\,10.42$\times$, eval collapses to 31.7). Blocking reward gradients forces co-adaptation through the observation decoder, reinforcing spatial bias.
    \item \textbf{Full-pathway SGS}: Block all non-KL gradients through \texttt{dt\_embedding}. Result: 0/3 across Walker and Cheetah. Even with only KL dynamics gradients updating the embedding, CIVO persists.
    \item \textbf{Obs-decoder-only SGS}: Block only observation decoder gradients (the identified causal channel for spatial bias). Result: 1/6 across Walker (4 seeds) and Cheetah (2 seeds). The single surviving Walker seed still suffers universal value collapse (value-to-eval ratio 0.056 at 215K). On Cheetah, NVE \emph{amplifies} to 21.9--30.2$\times$ (vs.\ untreated 8.3$\times$)---a compensatory amplification effect.
\end{enumerate}

The aggregate result: \textbf{1/17 seeds avoid CIVO}, and that seed still suffers critic calibration damage.
By contrast, embedding removal (008f + 024b) achieves 100\% CIVO prevention (5/5 seeds).

\textbf{Core finding}: Gradient magnitude is decoupled from causal importance.
The reward head contributes the largest gradients (38--49\%) but reward-only SGS fails worst.
The observation decoder contributes the smallest non-zero gradients (16--22\%) but targeted SGS also fails.
The failure is fundamental: backward-pass interventions cannot undo forward-pass representational damage.
The \texttt{dt\_embedding} shapes RSSM latent representations through the forward pass regardless of which loss components update it.

%% ---- 4.7 C8 ----
\subsection{Forward-Pass Mitigation: Temporal Conditioning Noise (C8)}
\label{sec:c8}

The backward-pass failure diagnosis (Section~\ref{sec:c7}) motivates forward-pass interventions that directly reduce the precision of temporal conditioning during imagination.
We propose \textbf{Temporal Conditioning Noise (TCN)}: during imagination rollouts, inject Gaussian noise $\epsilon \sim \mathcal{N}(0, \sigma^2)$ into the \texttt{dt\_embedding} output before it conditions the RSSM.
The noise prevents the critic from relying on precise (but inaccurate) temporal signals, disrupting the co-adaptation that underlies CIVO.

\paragraph{TCN $\sigma=0.2$, Walker Walk.}
NVE remains below 1.0$\times$ throughout training (peak 0.908 at 95K, terminal 0.095 at 215K).
Evaluation peaks at 821.1 (90\% of baseline), demonstrating that CIVO prevention does not require sacrificing performance.
This is the first intervention across the entire experimental campaign to achieve complete CIVO prevention with competitive performance.

\paragraph{TCN $\sigma=0.5$, Walker Walk (2 seeds).}
Excessive noise creates an oscillatory bistable regime: NVE peaks at 1.6--4.0$\times$, evaluation plateaus at $\sim$365 (40\% baseline).
Useful temporal information is destroyed along with the pathological signal.

\paragraph{TCN $\sigma=0.2$, Cheetah Run.}
Insufficient dose: NVE\,=\,1.903 at 55K (just exceeds 1.5$\times$ threshold).
Cheetah's untreated CIVO severity is 10--35$\times$ that of Walker (NVE 8--70$\times$ vs.\ 2--4$\times$), requiring a proportionally larger noise dose.

\paragraph{TCN $\sigma=0.5$, Cheetah Run.}
NVE reduced from untreated 8--108$\times$ to 1.5--4$\times$---a substantial reduction though not complete prevention.

\paragraph{Dose-response relationship.}
The required noise magnitude $\sigma$ correlates positively with untreated CIVO severity:
Walker (untreated NVE 2--4$\times$) is treated by $\sigma=0.2$;
Cheetah (untreated NVE 8--70$\times$) requires $\sigma=0.5$.
This provides mechanistic insight: the noise must be sufficient to degrade temporal conditioning precision below the threshold at which the critic can exploit RSSM fidelity gaps.

We also tested \textbf{RAVE (Reality-Anchored Value Estimation)}, which caps imagined values using an EMA of real environment values.
RAVE fails because the EMA anchor itself becomes polluted during CIVO onset---early overestimated values shift the anchor upward, rendering the ceiling ineffective.

An \textbf{uncertainty-gated imagination truncation} pilot (truncating imagination when RSSM prediction uncertainty exceeds a threshold) achieves partial reduction (NVE from 2--4$\times$ to $\sim$1.8$\times$, eval 380 at 105K) but is less effective than TCN.
This independently confirms that forward-pass interventions are the correct layer.

%% ---- 4.8 Cross-env ----
\subsection{Cross-Environment Generalization}
\label{sec:cross_env}

CIVO generalizes across all 4 tested environments, including the non-locomotion task Cartpole Swingup.

\paragraph{Cheetah Run: amplified severity.}
Cheetah's more complex dynamics (6-DOF vs.\ Walker's 2-DOF) and steeper reward landscape amplify CIVO: NVE reaches 8--70$\times$ (vs.\ Walker 2--4$\times$) with 100\% collapse rate (3/3 seeds).

\begin{table}[h]
\centering
\caption{CIVO severity scales with environment complexity.}
\label{tab:severity_cross_env}
\small
\begin{tabular}{@{}lccc@{}}
\toprule
Environment & NVE Range & Collapse Rate & Typical Onset \\
\midrule
Walker Walk      & 2--4$\times$   & $\sim$50\% (3/6)  & 35--55K \\
Cartpole Swingup & 1.9--4.5$\times$ & $\sim$67\% (2/3) & 45--55K \\
Hopper Hop       & 2.8--7.1$\times$ & $\sim$67\% (2/3) & 45--60K \\
Cheetah Run      & 8--70$\times$  & 100\% (3/3)        & 45--65K \\
\bottomrule
\end{tabular}
\end{table}

\paragraph{Cartpole Swingup: CIVO in non-locomotion tasks.}
All 3 seeds with trained \texttt{dt\_embedding} exceed NVE\,1.5$\times$ (range 1.85--4.5$\times$), with 2/3 collapsing.
This confirms the mechanism is architecture-level (RSSM + \texttt{dt\_embedding}), not task-specific.

\begin{table}[h]
\centering
\caption{Cartpole Swingup: trained dt\_embedding triggers CIVO in non-locomotion tasks.}
\label{tab:cartpole_civo}
\small
\begin{tabular}{@{}lcccc@{}}
\toprule
Seed & Eval Peak & Eval@105K & NVE Peak & Collapse \\
\midrule
s0 & 520 & 95  & 3.20 & Yes \\
s1 & 610 & 580 & 1.85 & No (self-corrected) \\
s2 & 480 & 70  & 4.50 & Yes \\
\bottomrule
\end{tabular}
\end{table}

\paragraph{NVE 1.5$\times$ diagnostic threshold.}
Across all 4 environments, the NVE\,$\geq$\,1.5$\times$ threshold at 35K+ steps achieves 26/26 correct classifications (Table~\ref{tab:nve_threshold_4env}).

\begin{table}[h]
\centering
\caption{NVE\,$\geq$\,1.5$\times$ threshold: cross-environment classification (trained dt\_emb exceeds threshold vs.\ untrained/absent stays below). AUROC\,=\,0.969.}
\label{tab:nve_threshold_4env}
\small
\begin{tabular}{@{}lccc@{}}
\toprule
Environment & Trained $>$1.5$\times$ & Untrained $<$1.5$\times$ & Total Correct \\
\midrule
Walker Walk      & 5/5 & 4/4 & 9/9 \\
Cheetah Run      & 3/3 & 2/2 & 5/5 \\
Cartpole Swingup & 3/3 & 3/3 & 6/6 \\
Hopper Hop       & 3/3 & 3/3 & 6/6 \\
\midrule
\textbf{Total}   & \textbf{14/14} & \textbf{12/12} & \textbf{26/26} \\
\bottomrule
\end{tabular}
\end{table}

Threshold sensitivity analysis (Table~\ref{tab:threshold_sensitivity}) confirms 1.5$\times$ as the optimal operating point (F1\,=\,0.95), balancing false positives (transient spikes at lower thresholds) against false negatives (missed borderline cases at higher thresholds).

\begin{table}[h]
\centering
\caption{NVE threshold sensitivity analysis.}
\label{tab:threshold_sensitivity}
\small
\begin{tabular}{@{}lcccc@{}}
\toprule
Threshold & Accuracy & FP & FN & F1 \\
\midrule
1.3$\times$ & 0.88 & 2 & 0 & 0.90 \\
\textbf{1.5$\times$} & \textbf{0.94} & \textbf{0} & \textbf{1} & \textbf{0.95} \\
1.7$\times$ & 0.91 & 0 & 2 & 0.91 \\
2.0$\times$ & 0.84 & 0 & 4 & 0.83 \\
\bottomrule
\end{tabular}
\end{table}

\paragraph{Cross-environment causal control.}
Removing \texttt{dt\_embedding} eliminates CIVO across all environments (0/9+ seeds exceed NVE 1.5$\times$), with consistent reward degradation:

\begin{table}[h]
\centering
\caption{No-dt-embedding causal control: CIVO eliminated across all environments.}
\label{tab:no_emb_cross_env}
\small
\begin{tabular}{@{}lccc@{}}
\toprule
Environment & Seeds Tested & NVE $>$1.5$\times$ & CIVO Rate \\
\midrule
Walker Walk      & 4+ & 0 & 0\% \\
Cheetah Run      & 2+ & 0 & 0\% \\
Cartpole Swingup & 3  & 0 & 0\% \\
\midrule
Total            & 9+ & 0 & 0\% \\
\bottomrule
\end{tabular}
\end{table}

%% ============================================================
%%  5. DISCUSSION
%% ============================================================
\section{Discussion}

\paragraph{The CIVO mechanism.}
CIVO arises from the interaction of two factors: (1) the RSSM's inherently lower prediction fidelity at $dt>1$ (obs MSE gap 2.04$\times$ even in the baseline), and (2) a trained \texttt{dt\_embedding} that enables the critic to exploit this gap through precise temporal conditioning.
The result is a positive feedback loop: the critic overestimates values at $dt>1$ $\to$ the actor shifts its policy toward these overestimated regions $\to$ distribution shift degrades RSSM predictions $\to$ further overestimation.

\paragraph{Why MPC succeeds where imagination fails.}
TAWM \cite{nhu2025tawm} demonstrates successful temporal conditioning in TD-MPC2's MPC framework.
The key difference: MPC re-plans from real environment states at each step, breaking the positive feedback loop.
In imagination-based actor-critic, the critic bootstraps from its own imagined predictions, enabling CIVO's self-reinforcing dynamics.
This architectural distinction---not temporal conditioning per se---determines whether variable-rate planning succeeds.

\paragraph{Forward-pass vs.\ backward-pass interventions.}
The comprehensive failure of gradient-level interventions (1/17 seeds across 3 SGS variants, Section~\ref{sec:c7}) establishes that CIVO is fundamentally a forward-pass phenomenon.
The \texttt{dt\_embedding} shapes RSSM latent representations through the forward pass; its influence persists regardless of gradient routing.
Compensatory amplification (blocking one gradient channel strengthens co-adaptation through others) further confirms this conclusion.
TCN succeeds precisely because it operates at the forward-pass level, degrading the conditioning precision that enables exploitation without modifying gradient flow.

\paragraph{Dose-response and calibration.}
The TCN dose-response---Walker ($\sigma=0.2$) vs.\ Cheetah ($\sigma=0.5$)---suggests that the required noise scales with the severity of the untreated CIVO, which in turn scales with RSSM fidelity gap magnitude.
This is consistent with the mechanistic model: noise must degrade temporal conditioning precision below the exploitation threshold, and larger fidelity gaps require more noise to achieve this.
An adaptive noise schedule that calibrates $\sigma$ to measured RSSM fidelity could eliminate the need for per-environment tuning.

\paragraph{Multiple vulnerability pathways.}
The 008e\_v3 results (Section~\ref{sec:c6}) reveal that trained \texttt{dt\_embedding} creates vulnerability through at least two pathways.
Pathway~A (classic CIVO) operates through value overestimation and is eliminated by $dt=1$ imagination.
Pathway~B (policy instability) operates through RSSM latent distortion from $dt>1$ training data and persists even with $dt=1$ imagination.
This dual-pathway structure implies that surface-level fixes (restricting imagination dt) are insufficient; the \texttt{dt\_embedding} itself must be either removed or regularized (as TCN does).

\paragraph{Limitations.}
All experiments use DreamerV3 with RSSM.
While the CIVO mechanism likely applies to other latent-dynamics models with temporal conditioning, we do not test this.
The TCN dose requires environment-specific tuning ($\sigma=0.2$ for Walker, $\sigma=0.5$ for Cheetah); adaptive calibration is left to future work.
The NVE diagnostic threshold (1.5$\times$) is validated on DMC tasks; applicability to higher-dimensional observation spaces or sparse-reward tasks is untested.

%% ============================================================
%%  6. CONCLUSION
%% ============================================================
\section{Conclusion}

We present a systematic diagnostic analysis of why temporal abstraction fails in DreamerV3's imagination-based actor-critic.
The root cause is \emph{conditioning-induced value overestimation} (CIVO): a trained \texttt{dt\_embedding} enables the critic to exploit RSSM fidelity gaps at $dt>1$, triggering a positive feedback loop that collapses the policy.
Our NVE\,$\geq$\,1.5$\times$ diagnostic achieves 26/26 classification across 4 environments.
The exhaustive failure of backward-pass (gradient) interventions across 17 seeds establishes that CIVO is forward-pass representational damage, not gradient contamination.
Temporal Conditioning Noise (TCN), a forward-pass intervention that injects calibrated noise into the \texttt{dt\_embedding} during imagination, successfully prevents CIVO in Walker ($\sigma=0.2$, 90\% baseline performance) and substantially reduces it in Cheetah ($\sigma=0.5$).

The actionable takeaway: variable-rate imagination in DreamerV3-class architectures requires either (a)~removing temporal conditioning from the world model, at significant performance cost, or (b)~regularizing temporal conditioning precision via forward-pass noise injection.
Hierarchical approaches that maintain separate dynamics per temporal scale (e.g., THICK \cite{kim2024thick}, Hieros \cite{robine2024hieros}) may avoid flat-conditioning failure entirely and represent a promising direction.

%% ============================================================
%%  REFERENCES
%% ============================================================
\bibliographystyle{plainnat}

\begin{thebibliography}{20}

\bibitem[Hafner et~al.(2023)]{hafner2023mastering}
D.~Hafner, J.~Pasukonis, J.~Ba, and T.~Lillicrap.
\newblock Mastering diverse domains through world models.
\newblock \emph{arXiv preprint arXiv:2301.04104}, 2023.

\bibitem[Hafner et~al.(2021)]{hafner2021mastering}
D.~Hafner, T.~Lillicrap, M.~Norouzi, and J.~Ba.
\newblock Mastering {A}tari with discrete world models.
\newblock In \emph{ICLR}, 2021.

\bibitem[Hafner et~al.(2020)]{hafner2020dream}
D.~Hafner, T.~Lillicrap, J.~Ba, and M.~Norouzi.
\newblock Dream to control: Learning behaviors by latent imagination.
\newblock In \emph{ICLR}, 2020.

\bibitem[Sutton et~al.(1999)]{sutton1999between}
R.~S. Sutton, D.~Precup, and S.~Singh.
\newblock Between {MDP}s and semi-{MDP}s: A framework for temporal abstraction in reinforcement learning.
\newblock \emph{Artificial Intelligence}, 112(1--2):181--211, 1999.

\bibitem[Precup(2000)]{precup2000temporal}
D.~Precup.
\newblock \emph{Temporal Abstraction in Reinforcement Learning}.
\newblock PhD thesis, University of Massachusetts Amherst, 2000.

\bibitem[Sharma et~al.(2017)]{sharma2017learning}
S.~Sharma, A.~S. Lakshminarayanan, and B.~Ravindran.
\newblock Learning to repeat: Fine grained action repetition for deep reinforcement learning.
\newblock In \emph{ICLR}, 2017.

\bibitem[Lakshminarayanan et~al.(2017)]{lakshminarayanan2017dynamic}
A.~S. Lakshminarayanan, S.~Sharma, and B.~Ravindran.
\newblock Dynamic action repetition for deep reinforcement learning.
\newblock In \emph{AAAI}, 2017.

\bibitem[Micheli et~al.(2023)]{micheli2023transformers}
V.~Micheli, E.~Alonso, and F.~Fleuret.
\newblock Transformers are sample-efficient world models.
\newblock In \emph{ICLR}, 2023.

\bibitem[Robine et~al.(2023)]{robine2023transformer}
J.~Robine, M.~H{\"o}ftmann, T.~Uelwer, and S.~Harmeling.
\newblock Transformer-based world models are happy with 100k interactions.
\newblock In \emph{ICLR}, 2023.

\bibitem[Nhu et~al.(2025)]{nhu2025tawm}
Q.~Nhu, T.~Lan, and B.~Boots.
\newblock {TAWM}: Temporally-aware world models for model-based reinforcement learning.
\newblock In \emph{ICML}, 2025.

\bibitem[Lu et~al.(2018)]{lu2018non}
T.~Lu, D.~Schuurmans, and C.~Boutilier.
\newblock Non-delusional {Q}-learning and value-iteration.
\newblock In \emph{NeurIPS}, 2018.

\bibitem[Janner et~al.(2019)]{janner2019trust}
M.~Janner, J.~Fu, M.~Zhang, and S.~Levine.
\newblock When to trust your model: Model-based policy optimization.
\newblock In \emph{NeurIPS}, 2019.

\bibitem[Tassa et~al.(2018)]{tassa2018deepmind}
Y.~Tassa, Y.~Doron, A.~Muldal, T.~Erez, Y.~Li, D.~de~Las~Casas, D.~Budden, A.~Abdolmaleki, J.~Merel, A.~Lefrancq, T.~Lillicrap, and M.~Riedmiller.
\newblock {D}eep{M}ind {C}ontrol {S}uite.
\newblock \emph{arXiv preprint arXiv:1801.00690}, 2018.

\bibitem[Agarwal et~al.(2021)]{agarwal2021deep}
R.~Agarwal, M.~Schwarzer, P.~S. Castro, A.~C. Courville, and M.~Bellemare.
\newblock Deep reinforcement learning at the edge of the statistical precipice.
\newblock In \emph{NeurIPS}, 2021.

\bibitem[Kim et~al.(2024)]{kim2024thick}
S.~Kim, S.~Lee, and G.~Lee.
\newblock {THICK}: Temporal hierarchical imagination for model-based control.
\newblock \emph{arXiv preprint}, 2024.

\bibitem[Robine et~al.(2024)]{robine2024hieros}
J.~Robine, M.~H{\"o}ftmann, T.~Uelwer, and S.~Harmeling.
\newblock {H}ieros: Hierarchical imagination on structured state space sequence world models.
\newblock In \emph{ICML}, 2024.

\bibitem[Henderson et~al.(2018)]{henderson2018deep}
P.~Henderson, R.~Islam, P.~Bachman, J.~Pineau, D.~Precup, and D.~Meger.
\newblock Deep reinforcement learning that matters.
\newblock In \emph{AAAI}, 2018.

\end{thebibliography}

\end{document}
